\begin{lemma}[Characteristic function of the Gaussian measure $\gamma_n$]
Fix $n \in \mathbb{N}_0$ and let $\gamma_n$ be the Gaussian measure of
\noderef{GaussMeasure} on $\mathbb{R}^n$, with density
$\pi^{-n/2} e^{-|x|^2}$ with respect to Lebesgue measure.  Then for every
$t \in \mathbb{R}^n$ the characteristic function of $\gamma_n$, namely
$\widehat{\gamma_n}(t) = \int_{\mathbb{R}^n} e^{\,i \langle x, t\rangle}
  \, d\gamma_n(x)$, is given by
\[
  \int_{\mathbb{R}^n} e^{\,i \langle x, t\rangle} \, d\gamma_n(x)
    = e^{-|t|^2/4} .
\]
Note the denominator $4$ rather than the $2$ of the standard normal law: with the
normalization $d\gamma_n = \pi^{-n/2} e^{-|x|^2}\,dx$ of \noderef{GaussMeasure}
each coordinate has variance $\tfrac12$ rather than $1$, and the characteristic
function of a centred Gaussian of covariance $\sigma^2 I$ is
$e^{-\sigma^2|t|^2/2}$, which at $\sigma^2 = \tfrac12$ is $e^{-|t|^2/4}$.
\end{lemma}
\begin{proof}
Fix $t \in \mathbb{R}^n$.

Step 1 (unfolding the characteristic function and the density).  By definition
the characteristic function of a measure $\mu$ on an inner-product space at $t$
is $\int e^{\,i\langle x,t\rangle}\,d\mu(x)$, so the quantity to compute is
$\int_{\mathbb{R}^n} e^{\,i\langle x,t\rangle}\,d\gamma_n(x)$.  By
\noderef{GaussMeasure}, $\gamma_n$ is Lebesgue measure weighted by the density
\[
  \rho(x) = \pi^{-n/2} e^{-|x|^2} ,
\]
which is measurable and everywhere finite (indeed $0 < \rho(x) \le \pi^{-n/2}$).
Integration against a weighted measure is integration of the weighted integrand,
so
\[
  \int_{\mathbb{R}^n} e^{\,i\langle x,t\rangle} \, d\gamma_n(x)
    = \int_{\mathbb{R}^n} \rho(x)\, e^{\,i\langle x,t\rangle} \, dx ,
\]
the right-hand integral being with respect to Lebesgue measure on
$\mathbb{R}^n$.

Step 2 (rewriting the integrand).  For every $x \in \mathbb{R}^n$, using
$e^{a}e^{b} = e^{a+b}$, the symmetry $\langle x,t\rangle = \langle t,x\rangle$ of
the real inner product, and the factorization
$\rho(x) = \pi^{-n/2}\cdot e^{-|x|^2}$,
\[
  \rho(x)\, e^{\,i\langle x,t\rangle}
    = \pi^{-n/2}\, e^{-|x|^2}\, e^{\,i \langle t,x\rangle}
    = \pi^{-n/2}\, e^{\,-|x|^2 + i\langle t, x\rangle} .
\]
Since the two integrands agree at every point, their Lebesgue integrals agree,
and the constant $\pi^{-n/2}$ may be taken outside the integral by linearity:
\[
  \int_{\mathbb{R}^n} e^{\,i\langle x,t\rangle}\, d\gamma_n(x)
    = \pi^{-n/2} \int_{\mathbb{R}^n}
        e^{\,-|x|^2 + i \langle t, x\rangle} \, dx .
\]

Step 3 (the Gaussian Fourier integral).  The Gaussian Fourier transform in the
form valid on a real inner-product space $V$ of finite dimension $d$ states that
for complex $b$ with $\operatorname{Re} b > 0$, complex $c$ and $w \in V$,
\[
  \int_V e^{\,-b|x|^2 + c\,\langle w, x\rangle} \, dx
    = \Bigl(\frac{\pi}{b}\Bigr)^{d/2}
      \exp\Bigl(\frac{c^2 |w|^2}{4b}\Bigr) .
\]
Apply this with $V = \mathbb{R}^n$, so $d = \dim_{\mathbb{R}} V = n$, and with
$b = 1$ (whose real part is $1 > 0$), $c = i$ and $w = t$.  Since $i^2 = -1$,
\[
  \int_{\mathbb{R}^n} e^{\,-|x|^2 + i\langle t, x\rangle} \, dx
    = \pi^{n/2} \exp\Bigl(\frac{i^2 |t|^2}{4}\Bigr)
    = \pi^{n/2}\, e^{-|t|^2/4} .
\]

Step 4 (conclusion).  Substituting Step 3 into Step 2 and cancelling the
powers of $\pi$, which is legitimate because $\pi \neq 0$ and
$\pi^{-n/2}\pi^{n/2} = \pi^{-n/2 + n/2} = \pi^0 = 1$,
\[
  \int_{\mathbb{R}^n} e^{\,i\langle x,t\rangle} \, d\gamma_n(x)
    = \pi^{-n/2}\cdot \pi^{n/2}\, e^{-|t|^2/4}
    = e^{-|t|^2/4} ,
\]
which is the assertion.  As a check on the normalization, $t = 0$ gives
$\gamma_n(\mathbb{R}^n) = 1$, and the exponent $-|t|^2/4$ is the one recorded in
the statement.
\end{proof}
